\documentclass[10pt,a4paper]{amsart}
\usepackage{amsmath,amssymb,amsthm,graphicx}
% Retained minimal preamble for the Stacks Project chapter varieties.tex.
% Verbatim excerpt of the chapter preamble: Xy-pic setup (native preamble lines 14-18), the native theorem
% declarations with the mathematical macros (native preamble lines 154-217), and the project own no-op
% definitions of the reference/slogan/history display environments (native wrapper.tex lines 1-3) that the
% counted statement uses for its slogan. Copied mechanically; nothing added, removed or redefined.
% For commutative diagrams we use Xy-pic
\usepackage[all]{xy}

% We use 2cell for 2-commutative diagrams.
\xyoption{2cell}
% Theorem environments.
%
\theoremstyle{plain}
\newtheorem{theorem}[subsection]{Theorem}
\newtheorem{proposition}[subsection]{Proposition}
\newtheorem{lemma}[subsection]{Lemma}

\theoremstyle{definition}
\newtheorem{definition}[subsection]{Definition}
\newtheorem{example}[subsection]{Example}
\newtheorem{exercise}[subsection]{Exercise}
\newtheorem{situation}[subsection]{Situation}

\theoremstyle{remark}
\newtheorem{remark}[subsection]{Remark}
\newtheorem{remarks}[subsection]{Remarks}

\numberwithin{equation}{subsection}

% Macros
%
\def\lim{\mathop{\mathrm{lim}}\nolimits}
\def\colim{\mathop{\mathrm{colim}}\nolimits}
\def\Spec{\mathop{\mathrm{Spec}}}
\def\Hom{\mathop{\mathrm{Hom}}\nolimits}
\def\Ext{\mathop{\mathrm{Ext}}\nolimits}
\def\SheafHom{\mathop{\mathcal{H}\!\mathit{om}}\nolimits}
\def\SheafExt{\mathop{\mathcal{E}\!\mathit{xt}}\nolimits}
\def\Sch{\mathit{Sch}}
\def\Mor{\mathop{\mathrm{Mor}}\nolimits}
\def\Ob{\mathop{\mathrm{Ob}}\nolimits}
\def\Sh{\mathop{\mathit{Sh}}\nolimits}
\def\NL{\mathop{N\!L}\nolimits}
\def\CH{\mathop{\mathrm{CH}}\nolimits}
\def\proetale{{pro\text{-}\acute{e}tale}}
\def\etale{{\acute{e}tale}}
\def\QCoh{\mathit{QCoh}}
\def\Ker{\mathop{\mathrm{Ker}}}
\def\Im{\mathop{\mathrm{Im}}}
\def\Coker{\mathop{\mathrm{Coker}}}
\def\Coim{\mathop{\mathrm{Coim}}}

% Boxtimes
%
\DeclareMathSymbol{\boxtimes}{\mathbin}{AMSa}{"02}

%
% Macros for moduli stacks/spaces
%
\def\QCohstack{\mathcal{QC}\!\mathit{oh}}
\def\Cohstack{\mathcal{C}\!\mathit{oh}}
\def\Spacesstack{\mathcal{S}\!\mathit{paces}}
\def\Quotfunctor{\mathrm{Quot}}
\def\Hilbfunctor{\mathrm{Hilb}}
\def\Curvesstack{\mathcal{C}\!\mathit{urves}}
\def\Polarizedstack{\mathcal{P}\!\mathit{olarized}}
\def\Complexesstack{\mathcal{C}\!\mathit{omplexes}}
% \Pic is the operator that assigns to X its picard group, usage \Pic(X)
% \Picardstack_{X/B} denotes the Picard stack of X over B
% \Picardfunctor_{X/B} denotes the Picard functor of X over B
\def\Pic{\mathop{\mathrm{Pic}}\nolimits}
\def\Picardstack{\mathcal{P}\!\mathit{ic}}
\def\Picardfunctor{\mathrm{Pic}}
\def\Deformationcategory{\mathcal{D}\!\mathit{ef}}
\newenvironment{reference}{}{}
\newenvironment{slogan}{}{}
\newenvironment{history}{}{}
\usepackage{cleveref}
\crefname{originalreference}{original reference}{original references}
\Crefname{originalreference}{Original reference}{Original references}
\makeatletter
\newlabel{morphisms-lemma-scheme-over-field-universally-open}{{0383}{source}{}{sourcepointer-1}{}}
\newlabel{morphisms-lemma-scheme-over-field-universally-open@cref}{{[originalreference][0][]0383}{[1][0][]source}{}{}{}}
\makeatother
\crefname{originalreference}{original reference}{original references}
\Crefname{originalreference}{Original reference}{Original references}
\makeatletter
\newlabel{morphisms-lemma-integral-universally-closed}{{01WM}{source}{}{sourcepointer-2}{}}
\newlabel{morphisms-lemma-integral-universally-closed@cref}{{[originalreference][0][]01WM}{[1][0][]source}{}{}{}}
\makeatother
\crefname{originallemma}{original lemma}{original lemmas}
\Crefname{originallemma}{Original lemma}{Original lemmas}
\makeatletter
\newlabel{lemma-galois-action-irreducible-components-locally-finite-type}{{04KZ}{source}{}{sourcepointer-3}{}}
\newlabel{lemma-galois-action-irreducible-components-locally-finite-type@cref}{{[originallemma][0][]04KZ}{[1][0][]source}{}{}{}}
\makeatother
\crefname{originallemma}{original lemma}{original lemmas}
\Crefname{originallemma}{Original lemma}{Original lemmas}
\makeatletter
\newlabel{lemma-separably-closed-field-irreducible-components}{{038H}{source}{}{sourcepointer-4}{}}
\newlabel{lemma-separably-closed-field-irreducible-components@cref}{{[originallemma][0][]038H}{[1][0][]source}{}{}{}}
\makeatother
\crefname{originallemma}{original lemma}{original lemmas}
\Crefname{originallemma}{Original lemma}{Original lemmas}
\makeatletter
\newlabel{lemma-units-general-algebraically-closed}{{04L6}{source}{}{sourcepointer-5}{}}
\newlabel{lemma-units-general-algebraically-closed@cref}{{[originallemma][0][]04L6}{[1][0][]source}{}{}{}}
\makeatother
\crefname{originallemma}{original lemma}{original lemmas}
\Crefname{originallemma}{Original lemma}{Original lemmas}
\makeatletter
\newlabel{lemma-integral-closure-ground-field}{{04MI}{source}{}{sourcepointer-6}{}}
\newlabel{lemma-integral-closure-ground-field@cref}{{[originallemma][0][]04MI}{[1][0][]source}{}{}{}}
\makeatother
\title{Finite generation of the units modulo constants under any one of three additional hypotheses}
\author{The Stacks Project Authors}
\date{stacks-project:varieties.tex:local-snapshot-20261011; GNU FDL 1.2 or later}
\begin{document}
\maketitle
\noindent\textit{Editorial note.} All delivered mathematical text is the author's own, copied verbatim from the pinned Stacks Project chapter file varieties.tex; no statement or proof step was rewritten, reordered or removed.\par\medskip
\noindent\textit{Editorial note.} The counted span cites six cross-chapter labels, each recorded in assembly.external\_references with its canonical tag as a navigation pointer.\par\medskip
\noindent\textit{Editorial note.} The proof cites no bibliography item inside the counted span.\par\medskip
\noindent\textit{Allowed standard prerequisites (editorial).}\begin{itemize}
\item Finite generation of units for the algebraically closed case and integral closure of the ground field. Stacks Project, Varieties, Lemma lemma-units-general-algebraically-closed (tag 09NM) and Lemma lemma-integral-closure-ground-field (tag 09NN). The first is applied to the geometric component and the second identifies the integral closure with a field.
\item Irreducible components of a base change along a separably closed field. Stacks Project, Varieties, Lemma lemma-separably-closed-field-irreducible-components (tag 09NV). Ensures that the geometric component has the same irreducible components, so the algebraically closed case applies.
\end{itemize}\medskip
\ifdefined\hypertarget\hypertarget{sourcepointer-1}{}\fi
\noindent\textit{Source reference.} Original reference 0383: sources/stacks-chow/tags.txt, source lines 4222--4222. Cross-chapter standard prerequisite cited by the author inside the counted span; it is used as a pointer and its own proof is not part of this item.\par\medskip
\ifdefined\hypertarget\hypertarget{sourcepointer-2}{}\fi
\noindent\textit{Source reference.} Original reference 01WM: sources/stacks-chow/tags.txt, source lines 2487--2487. Cross-chapter standard prerequisite cited by the author inside the counted span; it is used as a pointer and its own proof is not part of this item.\par\medskip
\ifdefined\hypertarget\hypertarget{sourcepointer-3}{}\fi
\noindent\textit{Source reference.} Original lemma 04KZ: sources/stacks-chow/tags.txt, source lines 5990--5990. Cross-chapter standard prerequisite cited by the author inside the counted span; it is used as a pointer and its own proof is not part of this item.\par\medskip
\ifdefined\hypertarget\hypertarget{sourcepointer-4}{}\fi
\noindent\textit{Source reference.} Original lemma 038H: sources/stacks-chow/tags.txt, source lines 4236--4236. Cross-chapter standard prerequisite cited by the author inside the counted span; it is used as a pointer and its own proof is not part of this item.\par\medskip
\ifdefined\hypertarget\hypertarget{sourcepointer-5}{}\fi
\noindent\textit{Source reference.} Original lemma 04L6: sources/stacks-chow/tags.txt, source lines 5997--5997. Cross-chapter standard prerequisite cited by the author inside the counted span; it is used as a pointer and its own proof is not part of this item.\par\medskip
\ifdefined\hypertarget\hypertarget{sourcepointer-6}{}\fi
\noindent\textit{Source reference.} Original lemma 04MI: sources/stacks-chow/tags.txt, source lines 6044--6044. Cross-chapter standard prerequisite cited by the author inside the counted span; it is used as a pointer and its own proof is not part of this item.\par\medskip
\begin{proposition}
\label{proposition-units-general}
Let $k$ be a field. Let $X$ be a scheme over $k$. Assume that $X$ is
locally of finite type over $k$, connected, reduced, and has finitely many
irreducible components. Then $\mathcal{O}(X)^*/k^*$ is a finitely generated
abelian group if in addition to the conditions above at least
one of the following conditions is satisfied:
\begin{enumerate}
\item the integral closure of $k$ in $\Gamma(X, \mathcal{O}_X)$ is $k$,
\item $X$ has a $k$-rational point, or
\item $X$ is geometrically integral.
\end{enumerate}
\end{proposition}

\begin{proof}
Let $\overline{k}$ be an algebraic closure of $k$.
Let $Y$ be a connected component of $(X_{\overline{k}})_{red}$.
Note that the canonical morphism $p : Y \to X$ is open (by
Morphisms, Lemma \ref{morphisms-lemma-scheme-over-field-universally-open})
and closed (by
Morphisms, Lemma \ref{morphisms-lemma-integral-universally-closed}).
Hence $p(Y) = X$ as $X$ was assumed connected. In particular, as
$X$ is reduced this implies $\mathcal{O}(X) \subset \mathcal{O}(Y)$. By
Lemma \ref{lemma-galois-action-irreducible-components-locally-finite-type}
we see that $Y$ has finitely many irreducible components (the field
$\overline{k}$ in the lemma is the separable algebraic closure, but
that this doesn't matter as the two base changes $X_{\overline{k}}$
have the same irreducible components by
Lemma \ref{lemma-separably-closed-field-irreducible-components}). Thus
Lemma \ref{lemma-units-general-algebraically-closed}
applies to $Y$. This implies that if $\mathcal{O}(X)^*/k^*$ is
not a finitely generated abelian group, then there exist elements
$f \in \mathcal{O}(X)$, $f \not \in k$ which map to an element of
$\overline{k}$ via the map $\mathcal{O}(X) \to \mathcal{O}(Y)$.
In this case $f$ is algebraic over $k$, hence
integral over $k$. Thus, if condition (1) holds, then this cannot happen.
To finish the proof we show that conditions (2) and (3) imply (1).

\medskip\noindent
Let $k \subset k' \subset \Gamma(X, \mathcal{O}_X)$ be the integral
closure of $k$ in $\Gamma(X, \mathcal{O}_X)$. By
Lemma \ref{lemma-integral-closure-ground-field}
we see that $k'$ is a field.
If $e : \Spec(k) \to X$ is a $k$-rational point, then
$e^\sharp : \Gamma(X, \mathcal{O}_X) \to k$ is a section to the
inclusion map $k \to \Gamma(X, \mathcal{O}_X)$. In particular the
restriction of $e^\sharp$ to $k'$ is a field map $k' \to k$ over $k$,
which clearly shows that (2) implies (1).

\medskip\noindent
If the integral closure $k'$ of $k$ in $\Gamma(X, \mathcal{O}_X)$
is not trivial, then we see that $X$ is either not geometrically connected
(if $k'/k$ is not purely inseparable) or that $X$ is not
geometrically reduced (if $k'/k$ is nontrivial purely inseparable).
Details omitted. Hence (3) implies (1).
\end{proof}

\end{document}
